\chapter{Lektionsplanung}

Jede Doppellektion dauert 90 Minuten. Die Doppellektionen 1 bis 7 folgen demselben Ablauf: 10 Minuten Einstiegsfrage, 5 Minuten Impuls, 60 Minuten Arbeitsblock, 15 Minuten Austausch. Die Doppellektion 8 besteht aus Präsentationen.

% Argumente: Nummer, frei, Thema, Einstiegsfrage, Impuls, Arbeitsblock, Austausch
\newcommand{\lplan}[7]{%
	\subsection*{DL #1: #3}
	\begin{center}
		\small
		\begin{tabular}{@{}p{.07\textwidth}p{.19\textwidth}p{.65\textwidth}@{}}
			\midrule
			10' & Einstiegsfrage & #4 \\
			5' & Impuls & #5 \\
			60' & Arbeitsblock & #6 \\
			15' & Austausch & #7 \\
			\midrule
		\end{tabular}
	\end{center}}

\lplan{1}{}{Zufall simulieren}
{Geld mit 80 und Casino: Antworten sammeln, keine Auflösung. Danach Würfelspiel im Raum.}
{Notebook öffnen, Zelle ausführen, Selbstkontrolle erklären.}
{Zu zweit \texttt{01\_Zufall.ipynb}: Roulette, Casino-Bilanz, Challenge Genauigkeit.}
{Weshalb verliert das Casino nie? Was bedeutet Durchschnitt für den Einzelfall?}

\lplan{2}{}{Anlageformen}
{Was besitzen Sie, wenn das Geld auf der Bank liegt? Kurze schriftliche Antworten.}
{Forderung, Beteiligung, Risiko.}
{Steckbrief einer Anlageform zu zweit (30'), dann \texttt{02\_Anlageformen.ipynb} (30').}
{Galerierundgang: Wer verliert in welcher Anlage zuerst?}

\lplan{3}{}{30 Jahre in die Zukunft}
{Zwei Personen mit gleichem Plan, zwölffacher Unterschied: War die eine klüger?}
{Was ein Modellparameter ist und weshalb er keine Prognose ist.}
{Zu zweit \texttt{03\_30\_Jahre.ipynb}: Verläufe, Perzentile, Challenges Median und Sparplan.}
{Wie würden Sie das Ergebnis einer Person ohne Mathematik erklären?}

\lplan{4}{}{Der Durchschnitt täuscht}
{Münzspiel mit der ganzen Gruppe: Wer spielt mit und weshalb?}
{Mittelwert und Median an einem Beispiel.}
{Zu zweit \texttt{04\_Durchschnitt.ipynb}: Spiel, bester Einsatz, Reihenfolge der Renditen.}
{Lässt sich die Frage \enquote{Geld mit 80} nun beantworten? Was fehlt noch?}

\lplan{5}{}{Diversifikation}
{500 Aktien und trotzdem grosser Verlust: Wie kann das sein?}
{Korrelation an einem Bild.}
{\texttt{05\_Diversifikation.ipynb} (35'), danach Projektidee wählen und zu zweit festhalten (25').}
{Projektideen in je einem Satz vorstellen.}

\subsection*{Zwischen DL 5 und DL 6}
Die Projektgruppen geben ihre Leitfrage ab.

\lplan{6}{}{Versicherung im Alltag}
{Wer zahlt die Rechnung der Kranken? Kurze Standpunkte, auch rechtlich.}
{Franchise, Selbstbehalt, Solidarität.}
{\texttt{06\_Franchise.ipynb} (35'), danach Projektarbeit in \texttt{08\_Projektvorlage.ipynb} (25').}
{Ist das Modellergebnis gerecht? Was sagt das Modell nicht?}

\lplan{7}{}{Wenn das Modell versagt}
{Alle 13'800 Jahre, und dann zweimal in zehn Jahren: Wer haftet?}
{Dicke Ränder an einer Grafik.}
{\texttt{07\_Modell.ipynb} (35'), danach Fallbeispiel in Gruppen oder Projektarbeit (25').}
{Kurzer Fallbericht: Annahme, Ereignis, Geschädigte. Auftrag für die Eigenarbeit.}

\subsection*{Zwischen DL 7 und DL 8}
Eigenarbeit am Projekt, rund 8 Stunden.

\subsection*{DL 8: Präsentationen}
\begin{center}
	\small
	\begin{tabular}{@{}p{.07\textwidth}p{.19\textwidth}p{.65\textwidth}@{}}
		\midrule
		70' & Präsentationen & Bis zu sieben Gruppen mit je 10 Minuten inklusive Rückfrage \\
		20' & Abschluss & Rückkehr zur Frage \enquote{Geld mit 80}, Rückmeldung zum Kurs \\
		\midrule
	\end{tabular}
\end{center}
